%% ============================================================================
%% Fichier  : annexes/Ann_C_templates.tex
%% Titre    : Papiers de Travail Types
%% Version  : 0.1 -- novembre 2026
%% ============================================================================

\chapter{Papiers de Travail Types}
\label{ann:templates}

Les modèles ci-dessous reprennent les encadrés «~Papier de travail~» des
chapitres. Ils sont à recopier dans le tableur ou le traitement de texte
de l'institution.

\section*{C.1 Ordre de mission}

\begin{tableaucours}{L{3.6cm}L{8.6cm}}{%
    \tableauHeader{Rubrique} &
    \tableauHeader{À remplir}
}
Fondement & Texte, article, alinéa \\
Objectif & Une question à laquelle le rapport répondra \\
Périmètre & Entité, applications, ce qui est exclu \\
Période & Exercice ou dates \\
Durée et équipe & Dates, noms, chef de mission \\
Livrables & Rapport provisoire, contradictoire, rapport définitif \\
\end{tableaucours}

\section*{C.2 Demande de pièces}

\begin{tableaucours}{cL{5.6cm}L{3.0cm}L{1.6cm}L{1.6cm}}{%
    \tableauHeader{N°} &
    \tableauHeader{Pièce (extraite devant l'auditeur~?)} &
    \tableauHeader{Destinataire} &
    \tableauHeader{Délai} &
    \tableauHeader{Reçue le}
}
1 & & & & \\
2 & & & & \\
\end{tableaucours}

\section*{C.3 Questionnaire de contrôle interne}

\begin{tableaucours}{cL{4.6cm}L{2.6cm}L{2.4cm}L{2.0cm}}{%
    \tableauHeader{N°} &
    \tableauHeader{Question (qui, quand, quelle trace~?)} &
    \tableauHeader{Réponse déclarée} &
    \tableauHeader{Répondant} &
    \tableauHeader{Pièce}
}
1 & & & & \\
2 & & & & \\
\end{tableaucours}

\section*{C.4 Feuille de test}

\begin{tableaucours}{L{3.6cm}L{8.6cm}}{%
    \tableauHeader{Rubrique} &
    \tableauHeader{À remplir}
}
Référence & Domaine, contrôle, ligne du programme \\
Critère & Énoncé, source, force \\
Population & Source, nombre, totaux, empreinte \\
Échantillon & Méthode, taille, liste \\
Résultats & Élément par élément, avec la pièce \\
Cote & \OUI{}~/ \NON{}~/ \NEK{}, et constatation \\
Visa & Auteur, date, chef de mission \\
\end{tableaucours}

\section*{C.5 Fiche de constat}

\begin{tableaucours}{L{3.0cm}L{9.2cm}}{%
    \tableauHeader{Élément} &
    \tableauHeader{À remplir}
}
Condition & Fait chiffré, daté, sourcé \\
Critère & Règle citée et sa force \\
Cause & Ce qui, dans le dispositif, explique l'écart \\
Effet & Risque réalisé ou possible, chiffré \\
Recommandation & Action, responsable, échéance \\
Classement & Déficience ou déficience majeure, et pourquoi \\
\end{tableaucours}

\section*{C.6 Grille de classement des déficiences}

\begin{tableaucours}{L{2.6cm}L{2.2cm}L{2.2cm}L{2.2cm}L{1.6cm}L{1.4cm}}{%
    \tableauHeader{Déficience} &
    \tableauHeader{Perte possible} &
    \tableauHeader{Probabilité} &
    \tableauHeader{Compensation prouvée} &
    \tableauHeader{Fraude} &
    \tableauHeader{Classe}
}
A & & & & & \\
B & & & & & \\
\end{tableaucours}

\section*{C.7 Plan d'action et suivi}

\begin{tableaucours}{cL{4.6cm}L{2.4cm}L{1.8cm}L{1.6cm}L{1.6cm}}{%
    \tableauHeader{N°} &
    \tableauHeader{Recommandation} &
    \tableauHeader{Responsable} &
    \tableauHeader{Échéance} &
    \tableauHeader{Statut} &
    \tableauHeader{Preuve}
}
R-01 & & & & & \\
R-02 & & & & & \\
\end{tableaucours}
